% Generated by scripts/assemble.py; corrections belong in source/corrections.json.
\examhead{170}{西南交通大学 2024 年数学分析试卷(回忆版)}{162}
% p163-b01; source PDF page 163
\blockanchor{p163-b01}
\question{1}已知 $\lim_{n\to \infty}a_n = a$ ，证明： $\lim_{n\to \infty}\frac{[na_n]}{n} = a,$ 其中[·]表示取整.\par
% p163-b02; source PDF page 163
\blockanchor{p163-b02}
\question{2}设 $f(x)$ 在 $(a, b)$ 上连续，且 $f(a + 0), f(b - 0)$ 存在，证明：$f(x)$ 在 $(a, b)$ 上一致连续。\par
% p163-b03; source PDF page 163
\blockanchor{p163-b03}
\question{3}设 $f(x)$ 在 $[a, b]$ 上连续, 在 $(a, b)$ 上可导, 证明: 存在 $\xi \in (a, b)$, 使得 $f'(\xi) + \lambda f(\xi) = 0$.\par
% p163-b04; source PDF page 163
\blockanchor{p163-b04}
\question{4}忘了.\par
% p163-b05; source PDF page 163
\blockanchor{p163-b05}
\question{5}设 $f(x)$ 在 $[a, b]$ 上有界, 判断 $f, |f|, f^2$ 之间的可积性关系.\par
% p163-b06; source PDF page 163
\blockanchor{p163-b06}
\question{6}讨论级数 $\sum_{n=1}^{\infty} \frac{\sin nx}{n^p}$ 的敛散性.\par
% p163-b07; source PDF page 163
\blockanchor{p163-b07}
\question{7}证明 $f(x, y) = \sqrt{|xy|}$ 在点 (0,0) 处连续，并求 $f_x(0,0), f_y(0,0)$，判断 $f(x, y)$ 在点 (0,0) 的可微性.\par
% p163-b08; source PDF page 163
\blockanchor{p163-b08}
\question{8}设 $f(x)$ 在 $(0, +\infty)$ 上连续, 当 $x \to 0^{+}$ 时, $f(x) \to +\infty$, 且 $\int_{0}^{1} f(x) \mathrm{d}x$ 收敛, 证明: $\lim_{x \to 0^{+}} x f(x) = 0$.\par
% p163-b09; source PDF page 163
\par\Needspace{6\baselineskip}
\blockanchor{p163-b09}
\question{9}设 $u = f(x, y)$, 其中 $x = r \cos \theta, y = r \sin \theta$, 且 $\frac{\partial u}{\partial x} = \frac{\partial v}{\partial y}, \frac{\partial u}{\partial y} = -\frac{\partial v}{\partial x}$, 证明:\par
\nopagebreak[4]
% p163-b10; source PDF page 163
\blockanchor{p163-b10}
\displayformula{\frac {\partial^ {2} u}{\partial x ^ {2}} = \frac {1}{r ^ {2}} \frac {\partial^ {2} v}{\partial y ^ {2}}, \frac {\partial^ {2} u}{\partial y ^ {2}} = - \frac {1}{r ^ {2}} \frac {\partial^ {2} v}{\partial x ^ {2}}.}
% p163-b11; source PDF page 163
\blockanchor{p163-b11}
\question{10}求 $y = mx^2, y = nx^2 (0 < m < m)$, $y = \alpha x, y = \beta x (0 < \alpha < \beta)$ 所围区域的面积.\par
% p163-b12; source PDF page 163
\blockanchor{p163-b12}
\question{11}设 $\Omega$ 是单连通区域, $L$ 为其边界, 在 $\Omega$ 内部或外部取一定点, 设其为原点, $\pmb{r}$ 为 $L$ 上的点到原点的向量, $\pmb{n}$ 为该点的单位切向量, 求 $\oint_{L} \frac{\cos(\pmb{n}, \pmb{r})}{|\pmb{r}|} \mathrm{d}s$.\par
% p163-b13; source PDF page 163
\blockanchor{p163-b13}
\question{12}设 $f(x)$ 在 $[a, b]$ 上无界, 证明: 存在 $x_0 \in (a, b)$, 使得对任意的 $\delta > 0$, 满足 $f(x)$ 在 $(x_0 - \delta, x_0 + \delta) \cap [a, b]$ 上无界.\par
